% GENERATED FILE. Edit canonical lessons, then run npm run generate:books.
% canonical-source-hash: fnv1a64:c644be235e2686d7
% canonical-lessons: ML-C357-pinnitu, ML-C357-kramena, ML-C357-atuvare, ML-C357-annumutal, ML-C357-potunnane

\chapter{Afterwards and Ever Since}
\label{ch:ml-afterwards-and-ever-since}
\hlchaptermodality{357}

\begin{chapteropening}
\textbf{By the end of this chapter:} \emph{I can say afterwards, gradually, until then, ever since and all of a sudden.}

Complete the last lesson of chapter 357: I can say afterwards, gradually, until then, ever since and all of a sudden.
\end{chapteropening}

\section[\texorpdfstring{\ml{പിന്നീട്} (pinnīṭŭ)}{പിന്നീട് (pinnīṭŭ)}]{\ml{പിന്നീട്} (pinnīṭŭ) — afterwards; later}

\label{lesson:ML-C357-pinnitu}



\begin{quote}
Before the new one: say the Malayalam for to learn by hearsay, then the Malayalam for to fabricate.
\end{quote}

\subsection*{You'll want to know: \ml{പിന്നീട്}}
\textbf{\ml{പിന്നീട്}} — \emph{pinnīṭŭ} — \textquotedblleft{}afterwards; later\textquotedblright{}.

\begin{quote}
\ml{ആദ്യം} \ml{അവൻ} \ml{ഒരു} \ml{ലേഖകൻ} \ml{ആയിരുന്നു}, \ml{പിന്നീട്} \ml{പത്രാധിപരായി}.
\end{quote}

\begin{quote}
\emph{ādyaṁ avan oru lēkhakan āyirunnu, pinnīṭŭ patrādhiparāyi.}
\end{quote}

\begin{quote}
At first he was a correspondent, later he became an editor.
\end{quote}

\ml{പിന്നെ} is then, next; \ml{പിന്നീട്} is later on, at a later time.

\begin{cousinweb}
\textbf{\ml{പിന്നീട്}} (\emph{pinnīṭŭ}) is native Dravidian, built on \ml{പിൻ} (\emph{pin}), \textquotedblleft{}behind, after\textquotedblright{} (DEDR 4205), the root of \ml{പിന്നെ}, \textquotedblleft{}then\textquotedblright{}. It looks further ahead in time than \ml{പിന്നെ} does.
\end{cousinweb}

\begin{grammarlens}[title={What you've built}]
One more word for telling what happened.
\end{grammarlens}

\subsection*{Your turn}
\begin{itemize}
\raggedright
  \item \emph{Say it:} \emph{pinnīṭŭ}
  \item \emph{Say it:} \emph{pinnīṭŭ}, once more
  \item \emph{Say it:} \emph{ādyaṁ avan oru lēkhakan āyirunnu, pinnīṭŭ patrādhiparāyi.} — At first he was a correspondent, later he became an editor.
  \item \emph{Recall:} say the Malayalam for to learn by hearsay, then the Malayalam for to fabricate, then say \emph{pinnīṭŭ} again
\end{itemize}

\begin{tcolorbox}[breakable,colback=teal!4,colframe=teal!35!black,title={Before you move on}]
What does \ml{പിന്നീട്} mean? (\textquotedblleft{}Afterwards; later\textquotedblright{}.) Say it once more.
\end{tcolorbox}
\section[\texorpdfstring{\ml{ക്രമേണ} (kramēṇa)}{ക്രമേണ (kramēṇa)}]{\ml{ക്രമേണ} (kramēṇa) — gradually}

\label{lesson:ML-C357-kramena}



\begin{quote}
Before the new one: say the Malayalam for to fabricate, then the Malayalam for afterwards.
\end{quote}

\subsection*{You'll want to know: \ml{ക്രമേണ}}
\textbf{\ml{ക്രമേണ}} — \emph{kramēṇa} — \textquotedblleft{}gradually\textquotedblright{}.

\begin{quote}
\ml{മഴ} \ml{ക്രമേണ} \ml{കുറഞ്ഞു}.
\end{quote}

\begin{quote}
\emph{maḻa kramēṇa kuṟaññu.}
\end{quote}

\begin{quote}
The rain gradually eased.
\end{quote}

A Sanskrit form kept whole: it never changes its ending.

\begin{cousinweb}
\textbf{\ml{ക്രമേണ}} (\emph{kramēṇa}) is borrowed whole from Sanskrit \emph{krameṇa}, \textquotedblleft{}in regular course, by degrees\textquotedblright{} (CDIAL 3577), a form of \emph{krama}, \textquotedblleft{}a step, a series\textquotedblright{}: things change step by step.
\end{cousinweb}

\begin{grammarlens}[title={What you've built}]
One more word for telling what happened.
\end{grammarlens}

\subsection*{Your turn}
\begin{itemize}
\raggedright
  \item \emph{Say it:} \emph{kramēṇa}
  \item \emph{Say it:} \emph{kramēṇa}, once more
  \item \emph{Say it:} \emph{maḻa kramēṇa kuṟaññu.} — The rain gradually eased.
  \item \emph{Recall:} say the Malayalam for to fabricate, then the Malayalam for afterwards, then say \emph{kramēṇa} again
\end{itemize}

\begin{tcolorbox}[breakable,colback=teal!4,colframe=teal!35!black,title={Before you move on}]
What does \ml{ക്രമേണ} mean? (\textquotedblleft{}Gradually\textquotedblright{}.) Say it once more.
\end{tcolorbox}
\section[\texorpdfstring{\ml{അതുവരെ} (atuvare)}{അതുവരെ (atuvare)}]{\ml{അതുവരെ} (atuvare) — until then}

\label{lesson:ML-C357-atuvare}



\begin{quote}
Before the new one: say the Malayalam for afterwards, then the Malayalam for gradually.
\end{quote}

\subsection*{You'll want to know: \ml{അതുവരെ}}
\textbf{\ml{അതുവരെ}} — \emph{atuvare} — \textquotedblleft{}until then\textquotedblright{}.

\begin{quote}
\ml{അച്ഛൻ} \ml{ആറ്} \ml{മണിക്ക്} \ml{വന്നു}, \ml{അതുവരെ} \ml{ഞങ്ങൾ} \ml{കാത്തിരുന്നു}.
\end{quote}

\begin{quote}
\emph{acchan āṟŭ maṇikkŭ vannu, atuvare ñaṅṅaḷ kāttirunnu.}
\end{quote}

\begin{quote}
Father came at six, and until then we waited.
\end{quote}

Built like \ml{ഇതുവരെ}, so far: \ml{അത്}, that, + \ml{വരെ}, up to.

\begin{cousinweb}
\textbf{\ml{അതുവരെ}} (\emph{atuvare}) is native Dravidian: \ml{അത്} (\emph{atŭ}), \textquotedblleft{}that\textquotedblright{}, and \ml{വരെ}, \textquotedblleft{}up to, until\textquotedblright{}. It mirrors \ml{ഇതുവരെ}, \textquotedblleft{}until now\textquotedblright{}, with the far \textquotedblleft{}that\textquotedblright{} (DEDR 1) in place of the near \textquotedblleft{}this\textquotedblright{}.
\end{cousinweb}

\begin{grammarlens}[title={What you've built}]
One more word for telling what happened.
\end{grammarlens}

\subsection*{Your turn}
\begin{itemize}
\raggedright
  \item \emph{Say it:} \emph{atuvare}
  \item \emph{Say it:} \emph{atuvare}, once more
  \item \emph{Say it:} \emph{acchan āṟŭ maṇikkŭ vannu, atuvare ñaṅṅaḷ kāttirunnu.} — Father came at six, and until then we waited.
  \item \emph{Recall:} say the Malayalam for afterwards, then the Malayalam for gradually, then say \emph{atuvare} again
\end{itemize}

\begin{tcolorbox}[breakable,colback=teal!4,colframe=teal!35!black,title={Before you move on}]
What does \ml{അതുവരെ} mean? (\textquotedblleft{}Until then\textquotedblright{}.) Say it once more.
\end{tcolorbox}
\section[\texorpdfstring{\ml{അന്നുമുതൽ} (annumutal)}{അന്നുമുതൽ (annumutal)}]{\ml{അന്നുമുതൽ} (annumutal) — ever since; from that day}

\label{lesson:ML-C357-annumutal}



\begin{quote}
Before the new one: say the Malayalam for gradually, then the Malayalam for until then.
\end{quote}

\subsection*{You'll want to know: \ml{അന്നുമുതൽ}}
\textbf{\ml{അന്നുമുതൽ}} — \emph{annumutal} — \textquotedblleft{}ever since; from that day\textquotedblright{}.

\begin{quote}
\ml{അന്നുമുതൽ} \ml{അവൾ} \ml{രാവിലെ} \ml{നടക്കുന്നു}.
\end{quote}

\begin{quote}
\emph{annumutal avaḷ rāvile naṭakkunnu.}
\end{quote}

\begin{quote}
Ever since that day she has walked in the mornings.
\end{quote}

\ml{അന്ന്}, that day, + \ml{മുതൽ}, from: the time a change began.

\begin{cousinweb}
\textbf{\ml{അന്നുമുതൽ}} (\emph{annumutal}) is native Dravidian: \ml{അന്ന്} (\emph{annŭ}), \textquotedblleft{}on that day\textquotedblright{}, and \ml{മുതൽ} (\emph{mutal}), \textquotedblleft{}from, beginning with\textquotedblright{}, which DEDR 4950 sets beside Tamil \emph{mutal}, \textquotedblleft{}beginning, first\textquotedblright{}.
\end{cousinweb}

\begin{grammarlens}[title={What you've built}]
One more word for telling what happened.
\end{grammarlens}

\subsection*{Your turn}
\begin{itemize}
\raggedright
  \item \emph{Say it:} \emph{annumutal}
  \item \emph{Say it:} \emph{annumutal}, once more
  \item \emph{Say it:} \emph{annumutal avaḷ rāvile naṭakkunnu.} — Ever since that day she has walked in the mornings.
  \item \emph{Recall:} say the Malayalam for gradually, then the Malayalam for until then, then say \emph{annumutal} again
\end{itemize}

\begin{tcolorbox}[breakable,colback=teal!4,colframe=teal!35!black,title={Before you move on}]
What does \ml{അന്നുമുതൽ} mean? (\textquotedblleft{}Ever since; from that day\textquotedblright{}.) Say it once more.
\end{tcolorbox}
\section[\texorpdfstring{\ml{പൊടുന്നനെ} (poṭunnane)}{പൊടുന്നനെ (poṭunnane)}]{\ml{പൊടുന്നനെ} (poṭunnane) — all of a sudden}

\label{lesson:ML-C357-potunnane}



\begin{quote}
Before the new one: say the Malayalam for until then, then the Malayalam for ever since.
\end{quote}

\subsection*{You'll want to know: \ml{പൊടുന്നനെ}}
\textbf{\ml{പൊടുന്നനെ}} — \emph{poṭunnane} — \textquotedblleft{}all of a sudden\textquotedblright{}.

\begin{quote}
\ml{പൊടുന്നനെ} \ml{വിളക്കുകൾ} \ml{എല്ലാം} \ml{കെട്ടു}.
\end{quote}

\begin{quote}
\emph{poṭunnane viḷakkukaḷ ellāṁ keṭṭu.}
\end{quote}

\begin{quote}
All of a sudden all the lamps went out.
\end{quote}

A story-telling word, stronger than \ml{പെട്ടെന്ന്}.

\begin{cousinweb}
\textbf{\ml{പൊടുന്നനെ}} (\emph{poṭunnane}) is native Dravidian: DEDR 4483 lists Malayalam \emph{poṭunnanē}, \textquotedblleft{}suddenly, in a moment\textquotedblright{}, beside a Tamil expression for suddenness.
\end{cousinweb}

\begin{grammarlens}[title={What you've built}]
One more word for telling what happened.
\end{grammarlens}

\subsection*{Your turn}
\begin{itemize}
\raggedright
  \item \emph{Say it:} \emph{poṭunnane}
  \item \emph{Say it:} \emph{poṭunnane}, once more
  \item \emph{Say it:} \emph{poṭunnane viḷakkukaḷ ellāṁ keṭṭu.} — All of a sudden all the lamps went out.
  \item \emph{Recall:} say the Malayalam for until then, then the Malayalam for ever since, then say \emph{poṭunnane} again
\end{itemize}

\begin{tcolorbox}[breakable,colback=teal!4,colframe=teal!35!black,title={Before you move on}]
What does \ml{പൊടുന്നനെ} mean? (\textquotedblleft{}All of a sudden\textquotedblright{}.) Say it once more.
\end{tcolorbox}
